\chapter{Per-Agent Real-Time Estimation and Control}
\label{ch:peragent}

Each agent carries two constrained optimisation-based solvers as its onboard
computational core: a nonlinear model-predictive \emph{controller} that computes
the vehicle's inputs, and a moving-horizon \emph{estimator} that reconstructs its
state and unmeasured disturbances. The two are a matched pair, built on the same
real-time-iteration and constraint-handling machinery, and both are designed for
microsecond-to-millisecond execution on an embedded companion computer. This
chapter states their problem formulations, the real-time pattern in which they
are executed, and the application programming interface through which the
distributed layers of \cref{ch:dmhe,ch:dmpc} use them.

\begin{center}
\fcolorbox{accent!55}{accent!5}{\begin{minipage}{0.94\linewidth}
\vspace{2pt}
\textbf{Scope of this edition.} The solvers' internal algorithms---their
constraint handling, linear algebra and factorisation strategy---are the subject
of separate ongoing research by the author, extending an MSc thesis on
sequential-linear-quadratic-programming MPC, and are not disclosed here. They
are treated throughout this report as a \emph{fixed, validated capability},
specified by the problem classes they solve (\cref{eq:ocp,eq:mhe}), the
real-time pattern they execute in (\cref{sec:rti}), the interface they are
reached through (\cref{sec:capi}) and the timing envelope measured in
\cref{ch:simulation}. Every distributed result in this report depends only on
those four things: the coordination layers never see solver internals, and any
solver meeting the same interface can be substituted---as the accompanying
reference implementation demonstrates.
\vspace{3pt}
\end{minipage}}
\end{center}

\section{The model-predictive controller}
\label{sec:mpc}

\paragraph{Optimal control problem.} At each sampling instant, with the freshly
estimated state $\est{\state}$ as parameter, the controller transcribes the
finite-horizon optimal control problem by multiple shooting over a prediction
horizon $N$:
\begin{subequations}
\label{eq:ocp}
\begin{align}
\min_{\substack{\state_0,\dots,\state_N\\ \ctrl_0,\dots,\ctrl_{N-1}}}\quad
& \sum_{k=0}^{N-1}\ell_k(\state_k,\ctrl_k)+\ell_N(\state_N)\\
\text{s.t.}\quad
& \state_0=\est{\state},\\
& \state_{k+1}=\vv{f}_d(\state_k,\ctrl_k), && k=0,\dots,N-1,\\
& \vv{c}(\state_k,\ctrl_k)\le\vv{0}, && k=0,\dots,N-1,\\
& \vv{c}_N(\state_N)\le\vv{0},
\end{align}
\end{subequations}
where $\vv{f}_d$ is the RK4 discretisation of the vehicle dynamics
(\cref{ch:hexacopter,ch:fixedwing}), $\ell_k=\tfrac12\norm{\state_k-\state_k^{r}}_Q^2
+\tfrac12\norm{\ctrl_k-\ctrl_k^{r}}_R^2$ and $\ell_N=\tfrac12\norm{\state_N-\state_N^{r}}_{Q_N}^2$
are tracking-quadratic stage and terminal costs with weights $Q,Q_N\succeq0$,
$R\succ0$, and $\vv{c},\vv{c}_N$ are the general nonlinear stage and terminal
inequality constraints. Crucially for the UAV application, $\vv{c}$ carries
\emph{both} input constraints (thrust and torque limits, control-surface
saturations) \emph{and} state constraints (attitude and rate envelopes,
minimum airspeed, collision-avoidance rows contributed by the coordination layer
of \cref{ch:dmpc}); pure state constraints receive no special treatment.

\paragraph{Constraint handling.} The inequality constraints in \cref{eq:ocp} are
imposed by the solver itself, on the full nonlinear problem. They are not
softened into penalty terms, relaxed into box bounds on the inputs, or enforced
by clipping the computed command---the applied input $\ctrl_0$ of
\cref{eq:rti-feedback} already carries the geometry of whichever constraints are
active at the current linearisation point. This matters for the distributed
layers: the coupling constraints of \cref{ch:dmpc} are pushed into $\vv{c}$ as
additional rows through the interface of \cref{sec:capi}, and are then handled
by the same mechanism, at the same cost per row, as the vehicle's own actuator
and envelope limits.

\section{Real-time iteration}
\label{sec:rti}

Both solvers execute in the real-time-iteration (RTI) pattern of Diehl
\etal \cite{diehl2005realtime,gros2020linear}: rather than iterating a
nonlinear program to convergence within one sampling interval, each sample
performs a \emph{single} iteration on a problem that has itself shifted by one
step, so that the optimisation converges along the closed-loop trajectory rather
than within a sample. The cycle splits into two phases of very different cost:

\begin{itemize}[nosep,leftmargin=1.4em]
  \item a \textbf{preparation} phase, which performs the expensive work that
        does \emph{not} depend on the measurement about to arrive---model
        linearisation, sensitivity generation, and the matrix-level linear
        algebra---and which runs in the interval between samples;
  \item a \textbf{feedback} phase, which consumes the arriving measurement and
        emits the applied input immediately.
\end{itemize}

Because the first-stage feedback gain $\mat{K}_0$ and feedforward
$\vv{\kappa}_0$ are already available when the measurement arrives, the
measurement-to-actuation path is a single affine evaluation,
\begin{equation}
\boxed{\;\ctrl_0=\bar{\ctrl}_0+\mat{K}_0\,(\est{\state}-\bar{\state}_0)+\vv{\kappa}_0\;}
\label{eq:rti-feedback}
\end{equation}
about the current prediction $(\bar{\state}_0,\bar{\ctrl}_0)$. The latency from
measurement to input is therefore that of one gain multiplication---tens of
nanoseconds on an embedded core (\cref{sec:sim-compute})---while the remainder of the
cycle completes within the sampling interval. The estimator runs the mirror
image of this split, with an instantaneous state output on the arrival of
$\vv{y}_t$; \cref{sec:sim-compute} measures both phases for the vehicle models of
this report.

\section{The moving-horizon estimator}
\label{sec:mhe}

The estimator is the dual of the controller: instead of predicting inputs over a
future horizon, it fits states and disturbances to a past window of measurements.

\paragraph{Estimation problem.} At instant $t$, given the measurement history
$\vv{y}_{t-M},\dots,\vv{y}_t$, the known inputs $\ctrl_{t-M},\dots,\ctrl_{t-1}$ and
an \emph{arrival} pair $(\vv{x}_a,\Lambda_a)$ summarising all information before
the window, the estimator solves the window nonlinear least-squares problem over
$M$ intervals in the variables $\vv{z}=(\state_0,\vv{w}_0,\state_1,\dots,\vv{w}_{M-1},\state_M)$
(node $k$ holds $\vv{y}_k:=\vv{y}_{t-M+k}$):
\begin{subequations}
\label{eq:mhe}
\begin{align}
\min_{\vv{z}}\quad
&\tfrac12\norm{\state_0-\vv{x}_a}_{\Lambda_a}^2
+\sum_{k=0}^{M}\tfrac12\norm{\vv{y}_k-\vv{h}(\state_k,\ctrl_k)}_V^2
+\sum_{k=0}^{M-1}\tfrac12\norm{\vv{w}_k}_W^2\\
\text{s.t.}\quad
&\state_{k+1}=\vv{F}(\state_k,\ctrl_k,\vv{w}_k), && k=0,\dots,M-1,\\
&\vv{c}(\state_k,\vv{w}_k)\le\vv{0}, && k=0,\dots,M-1,\\
&\vv{c}_M(\state_M)\le\vv{0},
\end{align}
\end{subequations}
where $\vv{w}_k\in\R^{n_w}$ is an estimated disturbance input (an unmeasured
force, load, bias walk or parameter drift; ``input constraints'' on the
estimation side thus bound these disturbances), $V\succeq0$ is the measurement
information (zero rows encode missing or multi-rate measurements) and $W\succ0$
the disturbance information. The arrival term is the estimation-side counterpart
of the terminal cost; the constraints again carry both state rows (state-of-charge
limits, non-negativity, physical envelopes) and disturbance rows.

\paragraph{Window recursion and the arrival cost.} Problem \cref{eq:mhe} is
treated with the same constraint handling and the same real-time-iteration
pattern as the controller, one iteration per measurement, with the
instantaneous output described in \cref{sec:rti}. As the window advances, the
information carried by the stage that leaves it is not discarded but folded into
the arrival pair $(\vv{x}_a,\Lambda_a)$ by exact marginalisation---the standard
information-form construction of the moving-horizon arrival cost
\cite{rao2003constrained,rawlings2017model}---so that the estimator retains
its full history at constant per-sample cost.

\paragraph{Kalman-filter reduction.} A property we use directly in the
distributed estimator of \cref{ch:dmhe} is that, for a \emph{linear}
model, quadratic costs and \emph{no active constraints}, the solution of
\cref{eq:mhe} coincides exactly with that of the information/covariance Kalman
filter (equivalently, the extended Kalman filter for a linearised model). The
moving horizon estimator is therefore a strict generalisation of the Kalman
filter that adds (i) hard state and disturbance constraints, (ii) window
smoothing, and (iii) robustness to model nonlinearity, degrading gracefully to
the familiar filter when those are absent. In \cref{ch:dmhe} the translational
estimation sub-problem is linear, and we exploit this reduction for a
numerically robust per-agent filter.

\section{The C application programming interface}
\label{sec:capi}

Both solvers are exposed to the coordination layer through a small, allocation-%
free C interface with statically sized problem data, so that the C++ layer of
\cref{ch:architecture} never needs to know how the solver is built. The
controller interface is
\begin{verbatim}
  typedef struct { int nx,nu,nc,ncN,N; double Ts; } MPCDims;
  typedef struct MPCProblem MPCProblem;

  MPCProblem* mpc_create (const MPCDims *d);
  int  mpc_set_weights   (MPCProblem *p, const double *Q,
                          const double *R, const double *QN);
  int  mpc_set_reference (MPCProblem *p, const double *xref,
                          const double *uref);      // formation slot
  int  mpc_set_constraint_bounds(MPCProblem *p, const double *c_ub);
  int  mpc_prepare  (MPCProblem *p);          // before the measurement
  int  mpc_feedback (MPCProblem *p, const double *xhat, double *u0);
  void mpc_destroy  (MPCProblem *p);
\end{verbatim}
with a symmetric interface for the estimator (\texttt{mhe\_create},
\texttt{mhe\_init}, \texttt{mhe\_prepare(u\_prev,u\_t)},
\texttt{mhe\_feedback(y\_t,xhat)}, \texttt{mhe\_fuse\_relative(z,Pj)},
\texttt{mhe\_destroy}). The distributed layers call \texttt{mpc\_set\_reference},
\texttt{mpc\_set\_constraint\_bounds} and \texttt{mhe\_fuse\_relative} to inject
the coupling terms derived in the next chapters---formation references,
collision-avoidance rows and relative measurements---without any change to the
solver core. The solver, in the words of the architecture note that motivated
this design, ``does not even need to know that it is being used by a many-agent
distributed system.''

Two properties of this boundary are worth stating explicitly, because the rest of
the report relies on them. First, it is \emph{substitutable}: the accompanying
implementation ships a dependency-free reference behind exactly this interface
(a finite-horizon LQR for the controller, an information filter for the
estimator---precisely the linear reductions of \cref{sec:mhe}), so the entire
distributed stack compiles, runs and reproduces its results with no proprietary
component present. Second, it is \emph{instance-wise}: all state lives in the
caller-owned problem objects, so one process hosts as many independent solver
instances as it has agents, which is what makes the in-the-loop study of
\cref{sec:sim-inloop} possible.
